\documentclass[11pt,a4paper]{article}

\usepackage{amsmath,amssymb,amsthm}
\usepackage{mathtools}
\usepackage[margin=1in]{geometry}
\usepackage{stmaryrd}
\usepackage[round]{natbib}

% Custom inference rule macro
\newcommand{\infer}[3][]{%
  \genfrac{}{}{0.4pt}{0}{#3}{#2}\;\textsc{\footnotesize #1}%
}

\newtheorem{definition}{Definition}[section]
\newtheorem{remark}{Remark}[section]
\newtheorem{theorem}{Theorem}[section]
\newtheorem{lemma}[theorem]{Lemma}
\newtheorem{example}[theorem]{Example}

\newcommand{\C}{\mathsf{C}}
\newcommand{\Type}{\mathsf{Type}}
\newcommand{\ctx}{\;\mathsf{ctx}}
\newcommand{\type}{\;\mathsf{type}}
\newcommand{\at}[1]{\,@\,#1}
\newcommand{\one}{\mathbf{1}}
\newcommand{\zero}{\mathbf{0}}
\newcommand{\negc}{\neg}

\title{CredTT: A Type Theory for Graded Provability\\[0.5em]
\large Conditioning without Implication}

\author{Probabilistic Foundations Project}
\date{\today}

\begin{document}

\maketitle

\begin{abstract}
We define a type theory where judgments carry \emph{credences} from a De Morgan algebra. The algebra has multiplication, complement, zero, one, and order---with disjunction derived via De Morgan duality. When $\C = \{\zero, \one\}$, we recover Martin-L\"of Type Theory. When $\C = {[0,1]}$, we get probabilistic typing with conditioning via the chain rule (not division). The key insight: types and logic \emph{emerge} from credence structure. We add conditioning as a primitive, distinguishing this from purely algebraic approaches.
\end{abstract}

\section*{Terminological Note}

We use \emph{credence} (degree of belief) rather than \emph{probability} because:
\begin{enumerate}
    \item \textbf{Subjective interpretation}: Credences are epistemic---they represent degrees of belief or confidence in propositions, not frequencies of events.
    \item \textbf{Avoiding measure-theoretic baggage}: ``Probability'' suggests Kolmogorov axioms, sigma-algebras, and measure theory. Our credences satisfy only the De Morgan algebra axioms.
    \item \textbf{Type-theoretic context}: In type theory, we ask ``how confident are we that this type is inhabited?'' rather than ``what is the frequency of this event?''
    \item \textbf{Philosophical tradition}: The term ``credence'' is standard in formal epistemology for rational degrees of belief.
\end{enumerate}

Throughout this paper, $c$ denotes a credence value, $\C$ is the credence algebra, and ``credence'' replaces ``weight'' or ``probability'' in all technical contexts.

\tableofcontents

%=============================================================================
\section{The Credence Algebra}
%=============================================================================

\subsection{Product De Morgan Structure}

We axiomatize a minimal structure: multiplication and complement. Disjunction is derived.

\begin{definition}[Credence Algebra]
$\C$ is a type with:
\begin{itemize}
    \item Constants: $\one : \C$ (certainty), $\zero : \C$ (impossibility)
    \item Multiplication: $(\cdot) : \C \to \C \to \C$ (conjunction)
    \item Complement: $\negc : \C \to \C$ (negation)
    \item Order: $(\leq) : \C \to \C \to \Type$
\end{itemize}
\end{definition}

\begin{definition}[Multiplication Axioms]
\begin{align}
c \cdot \one &= c \tag{unit-right} \\
\one \cdot c &= c \tag{unit-left} \\
c \cdot \zero &= \zero \tag{annihilate-right} \\
\zero \cdot c &= \zero \tag{annihilate-left} \\
(u \cdot v) \cdot c &= u \cdot (v \cdot c) \tag{assoc} \\
u \cdot v &= v \cdot u \tag{comm}
\end{align}
\end{definition}

\begin{definition}[Complement Axioms]
\begin{align}
\negc \zero &= \one \tag{neg-zero} \\
\negc \one &= \zero \tag{neg-one} \\
\negc \negc c &= c \tag{involution}
\end{align}
\end{definition}

\begin{definition}[Order Axioms]
\begin{align}
c &\leq c \tag{refl} \\
u \leq v \land v \leq c &\Rightarrow u \leq c \tag{trans} \\
u \leq v \land v \leq u &\Rightarrow u = v \tag{antisym} \\
\zero &\leq c \tag{zero-least} \\
c &\leq \one \tag{one-greatest}
\end{align}
\end{definition}

\subsection{Derived Operations}

\begin{definition}[De Morgan Disjunction]
\[
c \vee v := \negc(\negc c \cdot \negc v)
\]
\end{definition}

When $\C = {[0,1]}$ with $\negc c = 1 - c$, we get:
\[
c \vee v = 1 - (1-c)(1-v) = c + v - cv
\]
This is the \textbf{probabilistic sum} (inclusion-exclusion for independent events).

\begin{theorem}[De Morgan Laws]
\begin{align}
\negc(c \cdot v) &= \negc c \vee \negc v \\
\negc(c \vee v) &= \negc c \cdot \negc v
\end{align}
\end{theorem}

\begin{proof}
Direct calculation from definitions.
\end{proof}

\subsection{What We Have vs What We Don't}

\begin{center}
\begin{tabular}{l|l}
\textbf{Primitive} & \textbf{Derived} \\
\hline
$\cdot$ (multiplication/AND) & $\vee$ (De Morgan OR) \\
$\negc$ (complement/NOT) & $\to$ (implication, if needed) \\
$\zero, \one$ (constants) & \\
$\leq$ (order) &
\end{tabular}
\end{center}

\textbf{NOT primitive}: addition $+$, subtraction $-$, division $/$.

\subsection{Instances}

\begin{example}[Boolean: $\C = \{\zero, \one\}$]
$\one \cdot \one = \one$, else $\zero$. $\negc \zero = \one$, $\negc \one = \zero$.

This is classical Boolean logic.
\end{example}

\begin{example}[Credence: $\C = {[0,1]}$]
$c \cdot v = c \times v$ (ordinary multiplication). $\negc c = 1 - c$.

This is product fuzzy logic / probabilistic reasoning.
\end{example}

\begin{example}[Three-valued: $\C = \{\zero, ?, \one\}$]
$? \cdot ? = ?$, $\negc ? = ?$.

Kleene's strong three-valued logic.
\end{example}

%=============================================================================
\section{Judgments}
%=============================================================================

\begin{definition}[Judgment Forms]
\begin{enumerate}
    \item $\Gamma \ctx$ \hfill (valid context)
    \item $\Gamma \vdash A \type$ \hfill (valid type)
    \item $\Gamma \vdash a : A \at{c}$ \hfill (credence-annotated typing)
    \item $\Gamma \vdash a \equiv b : A \at{c}$ \hfill (credence-annotated equality)
\end{enumerate}
\end{definition}

\begin{remark}[Interpretation]
$\Gamma \vdash a : A \at{c}$ means: term $a$ inhabits type $A$ with credence $c$.
\begin{itemize}
    \item $c = \one$: $a$ definitely inhabits $A$
    \item $c = \zero$: $a$ does not inhabit $A$ (vacuously true)
    \item $\zero < c < \one$: $a$ partially/probably inhabits $A$
\end{itemize}
\end{remark}

%=============================================================================
\section{Structural Rules}
%=============================================================================

\begin{gather}
\infer[\textsc{Ctx-Empty}]
  {\cdot \ctx}
  {}
\qquad
\infer[\textsc{Ctx-Ext}]
  {\Gamma, x : A \ctx}
  {\Gamma \ctx \quad \Gamma \vdash A \type}
\\[1em]
\infer[\textsc{Var}]
  {\Gamma, x : A, \Delta \vdash x : A \at{\one}}
  {}
\\[1em]
\infer[\textsc{Weaken}]
  {\Gamma \vdash a : A \at{v}}
  {\Gamma \vdash a : A \at{c} \quad v \leq c}
\\[1em]
\infer[\textsc{Subst}]
  {\Gamma \vdash b[a/x] : B[a/x] \at{c \cdot v}}
  {\Gamma \vdash a : A \at{c} \quad \Gamma, x : A \vdash b : B \at{v}}
\end{gather}

%=============================================================================
\section{Type Formers}
%=============================================================================

\subsection{Functions (Dependent)}

\begin{gather}
\infer[\textsc{$\Pi$-Form}]
  {\Gamma \vdash (x : A) \to B \type}
  {\Gamma \vdash A \type \quad \Gamma, x : A \vdash B \type}
\\[1em]
\infer[\textsc{$\Pi$-Intro}]
  {\Gamma \vdash \lambda x. b : (x : A) \to B \at{c}}
  {\Gamma, x : A \vdash b : B \at{c}}
\\[1em]
\infer[\textsc{$\Pi$-Elim}]
  {\Gamma \vdash f\,a : B[a/x] \at{c \cdot v}}
  {\Gamma \vdash f : (x : A) \to B \at{c} \quad \Gamma \vdash a : A \at{v}}
\end{gather}

\textbf{Key rule}: credences multiply in elimination.

\subsection{Pairs (Dependent)}

\begin{gather}
\infer[\textsc{$\Sigma$-Form}]
  {\Gamma \vdash \Sigma(x : A). B \type}
  {\Gamma \vdash A \type \quad \Gamma, x : A \vdash B \type}
\\[1em]
\infer[\textsc{$\Sigma$-Intro}]
  {\Gamma \vdash (a, b) : \Sigma(x : A). B \at{c \cdot v}}
  {\Gamma \vdash a : A \at{c} \quad \Gamma \vdash b : B[a/x] \at{v}}
\\[1em]
\infer[\textsc{$\Sigma$-Elim$_1$}]
  {\Gamma \vdash \pi_1\,t : A \at{c}}
  {\Gamma \vdash t : \Sigma(x : A). B \at{c}}
\qquad
\infer[\textsc{$\Sigma$-Elim$_2$}]
  {\Gamma \vdash \pi_2\,t : B[\pi_1 t / x] \at{c}}
  {\Gamma \vdash t : \Sigma(x : A). B \at{c}}
\end{gather}

Pair credence = product of component credences (joint probability).

\subsection{Sums (Coproducts)}

\begin{gather}
\infer[\textsc{$+$-Form}]
  {\Gamma \vdash A + B \type}
  {\Gamma \vdash A \type \quad \Gamma \vdash B \type}
\\[1em]
\infer[\textsc{$+$-Intro$_L$}]
  {\Gamma \vdash \mathsf{inl}\,a : A + B \at{c}}
  {\Gamma \vdash a : A \at{c}}
\qquad
\infer[\textsc{$+$-Intro$_R$}]
  {\Gamma \vdash \mathsf{inr}\,b : A + B \at{c}}
  {\Gamma \vdash b : B \at{c}}
\\[1em]
\infer[\textsc{$+$-Elim}]
  {\Gamma \vdash \mathsf{case}\,e\,\mathsf{of}\,\mathsf{inl}\,x \Rightarrow l \mid \mathsf{inr}\,y \Rightarrow r : C \at{c \cdot v}}
  {\Gamma \vdash e : A + B \at{c} \quad \Gamma, x : A \vdash l : C \at{v} \quad \Gamma, y : B \vdash r : C \at{v}}
\end{gather}

\textbf{Key rule}: credences multiply in sum elimination. The scrutinee credence $c$ combines with the branch credence $v$ to give joint credence $c \cdot v$.

\subsection{Empty and Unit (Omitted)}

\textbf{Design decision}: Pure CredTT omits the Empty type ($\mathbf{0}$) and Unit type ($\mathbf{1}$).

\begin{itemize}
    \item \textbf{Impossibility}: Credence $\zero$ on any type captures ``does not inhabit.'' No special Empty type needed.
    \item \textbf{Certainty}: Credence $\one$ on any type captures ``definitely inhabits.'' No special Unit type needed.
\end{itemize}

This avoids the philosophical question ``what TYPE is impossibility/certainty?'' In MLTT, $\mathbf{0}$ and $\mathbf{1}$ are primitive types. In CredTT, impossibility and certainty are \emph{credences}, not types.

\begin{remark}[No Ex Falso Explosion]
Without $\mathbf{0}$-Elim (abort), there is no explosion rule. If $a : A \at{\zero}$, then $f\,a : B \at{c \cdot \zero} = \at{\zero}$---the zero credence propagates via multiplication, not via a special elimination rule.
\end{remark}

\begin{remark}[MLTT Compatibility]
When extending CredTT for MLTT compatibility, one may add:
\begin{gather*}
\infer[\textsc{$\mathbf{0}$-Form}]
  {\Gamma \vdash \mathbf{0} \type}
  {\Gamma \ctx}
\qquad
\infer[\textsc{$\mathbf{0}$-Elim}]
  {\Gamma \vdash \mathsf{abort}\,e : A \at{c}}
  {\Gamma \vdash e : \mathbf{0} \at{c}}
\\[1em]
\infer[\textsc{$\mathbf{1}$-Form}]
  {\Gamma \vdash \mathbf{1} \type}
  {\Gamma \ctx}
\qquad
\infer[\textsc{$\mathbf{1}$-Intro}]
  {\Gamma \vdash \star : \mathbf{1} \at{\one}}
  {\Gamma \ctx}
\end{gather*}
The Agda formalization omits these for philosophical purity.
\end{remark}

\subsection{Identity}

\begin{gather}
\infer[\textsc{Id-Form}]
  {\Gamma \vdash \mathsf{Id}_A(a, b) \type}
  {\Gamma \vdash a : A \at{\one} \quad \Gamma \vdash b : A \at{\one}}
\\[1em]
\infer[\textsc{Id-Intro}]
  {\Gamma \vdash \mathsf{refl} : \mathsf{Id}_A(a, a) \at{c}}
  {\Gamma \vdash a : A \at{c}}
\\[1em]
\infer[\textsc{Id-Elim}]
  {\Gamma \vdash J(M, d, p) : M[b, p] \at{c \cdot v}}
  {\Gamma \vdash p : \mathsf{Id}_A(a, b) \at{c} \quad \Gamma \vdash d : M[a, \mathsf{refl}] \at{v}}
\end{gather}

The J eliminator transports along an identity proof. Credences multiply: the proof credence $c$ combines with the base case credence $v$.

%=============================================================================
\section{Computation Rules}
%=============================================================================

The computation (beta) rules reduce eliminator-intro pairs. In each case, credence multiplication is preserved.

\subsection{$\Pi$-Computation}
\[
(\lambda x. b)\,a \equiv b[a/x] : B[a/x] \at{c \cdot v}
\]
when $\Gamma, x : A \vdash b : B \at{c}$ and $\Gamma \vdash a : A \at{v}$.

\subsection{$\Sigma$-Computation}
\begin{align}
\pi_1(a, b) &\equiv a : A \at{c \cdot v} \\
\pi_2(a, b) &\equiv b : B[a/x] \at{c \cdot v}
\end{align}
when $\Gamma \vdash a : A \at{c}$ and $\Gamma \vdash b : B[a/x] \at{v}$.

\subsection{$+$-Computation}
\begin{align}
\mathsf{case}\,(\mathsf{inl}\,a)\,l\,r &\equiv l[a/x] : C \at{c \cdot v} \\
\mathsf{case}\,(\mathsf{inr}\,b)\,l\,r &\equiv r[b/y] : C \at{c \cdot v}
\end{align}

\subsection{Id-Computation}
\[
J(M, d, \mathsf{refl}) \equiv d : M[a, \mathsf{refl}] \at{c \cdot v}
\]
when $\Gamma \vdash a : A \at{c}$ and $\Gamma \vdash d : M[a, \mathsf{refl}] \at{v}$.

%=============================================================================
\section{Negation}
%=============================================================================

There are two approaches to negation:

\subsection{Type-Theoretic Negation}

\[
\neg A := A \to \mathbf{0}
\]

A proof of $\neg A$ is a function from $A$ to the empty type.

\subsection{Credence Negation}

For a term $a : A \at{c}$, we can consider $\negc c$ as the ``complementary credence.''

In $\C = {[0,1]}$: if $P(A) = c$, then $P(\neg A) = 1 - c = \negc c$.

\begin{remark}[Relationship]
Type-theoretic negation ($A \to \mathbf{0}$) corresponds to credence negation in the following sense: if $a : A \at{c}$ with $c > \zero$, then there is no $f : A \to \mathbf{0} \at{\one}$.
\end{remark}

%=============================================================================
\section{Conditioning (Not Implication)}
%=============================================================================

\subsection{Why Not Material Implication}

Material implication $A \to B := \negc A \vee B$ has paradoxes:
\begin{itemize}
    \item $\mathsf{False} \to \mathsf{Anything} = \mathsf{True}$ (ex falso)
    \item $\mathsf{Anything} \to \mathsf{True} = \mathsf{True}$
\end{itemize}

We use \textbf{conditioning} instead of implication.

\subsection{Chain Rule}

The fundamental equation is:
\[
P(A, B) = P(A \mid B) \cdot P(B) = P(B \mid A) \cdot P(A)
\]

This is \textbf{multiplicative}---no division needed!

\begin{gather}
\infer[\textsc{Chain}]
  {\Gamma \vdash (a, b) : A \times B \at{u \cdot v}}
  {\Gamma \vdash a : A \at{u} \quad \Gamma \vdash b \mid a : B \at{v}}
\end{gather}

The conditional $b \mid a$ is \textbf{primitive}. We don't compute it by division; we assert it exists when the chain rule holds.

\subsection{Conditioning vs Implication}

\begin{center}
\begin{tabular}{l|l|l}
& \textbf{Implication} $A \to B$ & \textbf{Conditioning} $A \mid B$ \\
\hline
Formula & $\negc A \vee B$ & $P(A,B) / P(B)$ \\
Computable from & $A$ and $B$ alone & Joint $P(A,B)$ \\
False antecedent & ``True'' (lies) & Indeterminate (honest) \\
In logic & Yes & No \\
In probability & No & Yes \\
\end{tabular}
\end{center}

\subsection{Graded Ex Falso}

What about $P(A \mid B)$ when $P(B) = 0$?

In classical logic: $\bot \to A$ is always true (ex falso quodlibet).

In CredTT: when $b : B \at{\zero}$:
\[
(a, b) : A \times B \at{\zero \cdot c} = \at{\zero}
\]
The conditional $a \mid b$ can have \emph{any} credence---the joint is already zero.

\textbf{Key insight}: Ex falso is the $c = 0$ limit of a continuous phenomenon:
\begin{itemize}
    \item $c = \one$: Evidence fully constrains conclusions
    \item $c = \epsilon$: Evidence weakly constrains conclusions
    \item $c = \zero$: Evidence imposes no constraint (ex falso)
\end{itemize}

Classical logic has only the endpoints; CredTT has the entire spectrum.

%=============================================================================
\section{MLTT as the $\{\zero, \one\}$ Case}
%=============================================================================

\begin{theorem}[MLTT Embedding]
When $\C = \{\zero, \one\}$, CredTT reduces to MLTT:
\begin{itemize}
    \item $\Gamma \vdash a : A \at{\one}$ becomes $\Gamma \vdash a : A$
    \item $\Gamma \vdash a : A \at{\zero}$ is vacuously satisfied
    \item Credence operations: $\one \cdot \one = \one$, else $\zero$
    \item Complement: $\negc\zero = \one$, $\negc\one = \zero$
    \item De Morgan OR: $c \vee v = \negc(\negc c \cdot \negc v)$ gives Boolean OR
\end{itemize}
\end{theorem}

\begin{proof}
By inspection. All credence operations compute to Boolean.
\end{proof}

\subsection{Side-by-Side Comparison}

\begin{center}
\begin{tabular}{l|l|l}
\textbf{Rule} & \textbf{MLTT} & \textbf{CredTT} \\
\hline
Variable & $x : A \vdash x : A$ & $x : A \vdash x : A \at{\one}$ \\
Application & $f\,a : B[a/x]$ & $f\,a : B[a/x] \at{c \cdot v}$ \\
Pair & $(a, b) : \Sigma A\,B$ & $(a, b) : \Sigma A\,B \at{c \cdot v}$ \\
Projection & $\pi_1 t : A$ & $\pi_1 t : A \at{c}$ \\
Case & $\mathsf{case}\,e\,l\,r : C$ & $\mathsf{case}\,e\,l\,r : C \at{c \cdot v}$ \\
\end{tabular}
\end{center}

\textbf{Key observation}: In MLTT, all judgments are implicitly at credence $\one$. Since $\one \cdot \one = \one$, credence multiplication is invisible. CredTT makes this structure explicit and generalizes it to arbitrary credences.

%=============================================================================
\section{Contrast with Cost-Based Graded Types}
%=============================================================================

Graded type systems like Granule \citep{orchard2019quantitative} and quantitative type theory \citep{atkey2018syntax} use grades to track \emph{resource usage}. CredTT uses credences to track \emph{inhabitation certainty}. These are philosophically distinct:

\begin{center}
\begin{tabular}{l|l|l}
& \textbf{Cost-Based (Granule)} & \textbf{CredTT} \\
\hline
Algebraic structure & Semiring $(0, 1, +, \cdot)$ & De Morgan $(\zero, \one, \cdot, \negc)$ \\
Addition & Primitive (branch costs add) & Derived via De Morgan \\
Philosophy & Decorate MLTT with grades & Logic \emph{emerges} from grades \\
Implication & Material implication exists & Implication \emph{rejected} \\
Interpretation & Costs accumulate & Credences condition/constrain \\
MLTT relationship & MLTT is axiomatic & MLTT is \emph{derivative} \\
\end{tabular}
\end{center}

\textbf{Key distinction}: Cost-based systems ask ``how much resource does this computation use?'' CredTT asks ``how certain are we that this type is inhabited?''

In cost-based grading:
\begin{itemize}
    \item $f : A \to B @ 2$ means ``$f$ uses 2 units of resource''
    \item Costs add at branches: if one branch costs 3 and another costs 5, total cost is at most 5
    \item The semiring axiom $0 + x = x$ reflects ``zero cost contributes nothing''
\end{itemize}

In CredTT:
\begin{itemize}
    \item $f : A \to B \at{0.9}$ means ``$f$ produces $B$ with 90\% certainty''
    \item Credences multiply in composition: chaining 90\% and 80\% gives 72\%
    \item The De Morgan axiom $\zero \cdot c = \zero$ reflects ``impossible antecedent gives impossible conclusion''
\end{itemize}

\textbf{One-liner}: Costs measure \emph{computation performed}; credences measure \emph{truth surviving composition}.

%=============================================================================
\section{Related Work: A Unified View}
%=============================================================================

CredTT synthesizes ideas from several traditions. We review each and explain how CredTT unifies them.

\subsection{Pavelka-Style Fuzzy Logic}

Pavelka \citep{pavelka1979fuzzy} introduced \emph{graded provability} for \textbf{propositional} fuzzy logic: instead of asking ``is $\varphi$ provable?'', one asks ``to what degree is $\varphi$ provable?''

\begin{definition}[Pavelka Provability Degree]
\[
|\varphi|_T := \sup\{r \in [0,1] \mid T \vdash (\varphi, r)\}
\]
\end{definition}

Pavelka \textbf{proved} a completeness theorem for this propositional system, equating syntactic provability degrees with semantic truth degrees.

H\'ajek \citep{hajek1998metamathematics} extended this to \textbf{predicate} fuzzy logic in his \emph{Metamathematics of Fuzzy Logic}. His completeness theorems apply to first-order fuzzy logic, not type theory.

\textbf{Key distinction}: Pavelka/H\'ajek work in propositional and first-order logic respectively, and they \textbf{prove} their completeness results. CredTT works in dependent type theory, which is a different formal system. The completeness question for CredTT (Section 11, Open Problem 7) remains open.

\textbf{What CredTT adapts from Pavelka/H\'ajek}:
\begin{itemize}
    \item The concept of graded provability (credences on judgments)
    \item The idea that degrees should form an algebra
\end{itemize}

\textbf{What CredTT adds}:
\begin{itemize}
    \item Dependent types ($\Pi$, $\Sigma$, $\mathsf{Id}$)
    \item Credences on \emph{judgments}, not just propositions
    \item Conditioning as primitive (not residuated implication)
\end{itemize}

\textbf{What CredTT postulates} (analogous to what Pavelka/H\'ajek prove):
\begin{itemize}
    \item That G\"odel sentences attain credence $\frac{1}{2}$ (we postulate the fixed point; they prove completeness)
\end{itemize}

\subsection{Self-Reference in Fuzzy Logic}

H\'ajek, Paris, and Shepherdson \citep{hajek2000liar} analyzed the liar paradox in fuzzy logic:

\begin{theorem}[H\'ajek-Paris-Shepherdson]
In {\L}ukasiewicz logic with a truth predicate $T$, the liar sentence $\lambda$ satisfying $\lambda \leftrightarrow \neg T(\ulcorner \lambda \urcorner)$ has truth value $\frac{1}{2}$.
\end{theorem}

\textbf{Key point}: H\'ajek-Paris-Shepherdson \textbf{prove} this result by carefully constructing the truth predicate and showing the fixed point exists. Their proof uses the diagonal lemma and properties of {\L}ukasiewicz logic.

\textbf{CredTT's version}: We \textbf{postulate} that the G\"odel sentence $G \leftrightarrow \neg\mathsf{Prov}_\one(G)$ exists and has credence satisfying $c = \negc c$, hence $c = \frac{1}{2}$. This uses the same fixed-point \emph{phenomenon}, but:
\begin{itemize}
    \item H\'ajek-Paris-Shepherdson work with \emph{truth} in {\L}ukasiewicz logic
    \item CredTT works with \emph{provability credences} in type theory
    \item They \textbf{construct} $\lambda$; we \textbf{postulate} $G$'s existence
\end{itemize}

A full treatment would require proving the diagonal lemma for CredTT and constructing $G$ explicitly.

\subsection{Logics of Formal Inconsistency (LFIs)}

Da Costa \citep{dacosta1974paraconsistent} and Carnielli-Coniglio \citep{carnielli2016paraconsistent} developed logics where contradiction doesn't explode:

\begin{definition}[Gentle Explosion]
\[
A, \neg A, \circ A \vdash B \quad \text{but} \quad A, \neg A \not\vdash B
\]
where $\circ A$ means ``$A$ is consistent.''
\end{definition}

LFIs have a well-developed proof theory with the consistency operator $\circ$ as \textbf{primitive syntax}. The operator gates explosion: you need $\circ A$ (assertion that $A$ is consistent) to derive arbitrary conclusions from $A \land \neg A$.

\textbf{CredTT's approach}: CredTT has no consistency operator and no explosion rule. From $a : A \at{c}$ and $b : \neg A \at{v}$, we simply cannot derive arbitrary $B \at{\one}$---there is no rule that would produce this. This is not ``gated explosion'' but \emph{absence of explosion}.

\begin{center}
\begin{tabular}{l|l|l}
& \textbf{LFIs} & \textbf{CredTT} \\
\hline
Explosion & Gated by $\circ$ operator & Absent (no explosion rule) \\
Formal device & Primitive unary connective & Graded judgments \\
Gradation & Binary (consistent/not) & Continuous $[0,1]$ \\
Proof theory & Well-developed & Postulated analogs \\
\end{tabular}
\end{center}

The comparison is suggestive but limited: LFIs and CredTT address paraconsistency through different mechanisms.

\subsection{T-Norm Fuzzy Logics}

Most fuzzy logics use \emph{t-norms} for conjunction \citep{klement2000triangular}:

\begin{definition}[T-norm]
A t-norm $* : [0,1]^2 \to [0,1]$ satisfies:
\begin{enumerate}
    \item $x * 1 = x$ (identity)
    \item $x * y = y * x$ (commutativity)
    \item $(x * y) * z = x * (y * z)$ (associativity)
    \item $x \leq y \Rightarrow x * z \leq y * z$ (monotonicity)
\end{enumerate}
\end{definition}

Common t-norms: product ($x \cdot y$), minimum ($\min(x,y)$), Lukasiewicz ($\max(0, x+y-1)$).

\textbf{CredTT uses product}: $c \cdot v$ is ordinary multiplication. This matches probabilistic conjunction (independent events) and gives the chain rule.

\textbf{Key difference}: Standard fuzzy logics include a \emph{residuum} (implication) defined by:
\[
x \to y := \sup\{z \mid x * z \leq y\}
\]

CredTT \emph{rejects} this implication in favor of conditioning. The residuum has material implication paradoxes; conditioning doesn't.

\subsection{Graded Type Theory (Granule)}

Orchard et al. \citep{orchard2019quantitative} and Atkey \citep{atkey2018syntax} developed graded type systems for resource tracking:

\begin{center}
\begin{tabular}{l|l|l}
& \textbf{Granule/QTT} & \textbf{CredTT} \\
\hline
Algebraic structure & Semiring $(0,1,+,\cdot)$ & De Morgan $(\zero,\one,\cdot,\negc)$ \\
Addition & Primitive (costs add) & Derived (De Morgan) \\
Interpretation & Resource usage & Inhabitation certainty \\
\end{tabular}
\end{center}

\textbf{Technical distinction}: Granule's semiring includes addition as primitive, enabling resource branching: ``use resource $r$ in one branch OR resource $s$ in another.'' CredTT derives addition via De Morgan duality, which limits expressiveness for branching scenarios.

\textbf{Philosophical distinction}: Both systems treat grades as primitive. The difference is interpretive:
\begin{itemize}
    \item Granule: ``how much resource does this computation use?''
    \item CredTT: ``how certain is this type inhabited?''
\end{itemize}

In Granule, $f : A \to B @ 2$ means ``$f$ uses 2 units.'' In CredTT, $f : A \to B \at{0.9}$ means ``$f$ produces $B$ with 90\% certainty.''

\textbf{Note}: The claim that MLTT ``emerges'' from CredTT while Granule ``decorates'' MLTT is a matter of presentation. Both can be viewed as parameterized systems that specialize to MLTT. Granule's semiring can also be instantiated at $\{0,1\}$.

\subsection{De Morgan Algebras}

Our credence algebra is a \emph{De Morgan algebra} \citep{kalman1958lattices}:

\begin{definition}[De Morgan Algebra]
$(L, \land, \lor, \neg, 0, 1)$ with:
\begin{itemize}
    \item $(L, \land, \lor, 0, 1)$ is a bounded distributive lattice
    \item $\neg\neg x = x$ (involution)
    \item $\neg(x \land y) = \neg x \lor \neg y$ (De Morgan)
\end{itemize}
\end{definition}

CredTT uses a \emph{product De Morgan algebra}: conjunction is multiplication, disjunction derived via De Morgan. When $L = [0,1]$ with $\neg c = 1-c$:
\begin{align}
c \land v &= c \cdot v \\
c \lor v &= \negc(\negc c \cdot \negc v) = c + v - cv
\end{align}

\subsection{Markov Categories}

Fritz \citep{fritz2020synthetic} developed Markov categories for probabilistic reasoning:

\begin{definition}[Markov Category]
A symmetric monoidal category with:
\begin{itemize}
    \item Copy: $\Delta_A : A \to A \otimes A$
    \item Delete: $!_A : A \to I$
    \item Naturality conditions for stochastic maps
\end{itemize}
\end{definition}

CredTT's categorical semantics is a \emph{credence-weighted Markov category}: morphisms carry credences, composition multiplies credences.

\subsection{Synthesis: What CredTT Unifies}

\begin{center}
\begin{tabular}{l|l}
\textbf{From} & \textbf{CredTT takes} \\
\hline
Pavelka/Hajek & Graded provability \\
Hajek-Paris-Shepherdson & Self-reference $\to$ fixed point at $\frac{1}{2}$ \\
LFIs & Non-explosive contradiction \\
T-norm logics & Product conjunction \\
Granule/QTT & Dependent types + grades \\
De Morgan algebras & Minimal algebraic structure \\
Markov categories & Conditioning semantics \\
\end{tabular}
\end{center}

\textbf{What CredTT combines}: Dependent type theory with conditioning (not implication), using De Morgan algebra for credences. MLTT is the $\{\zero,\one\}$ specialization.

\textbf{Honest assessment of novelty}: The individual components (graded provability, fuzzy fixed points, graded types, De Morgan algebras) are well-established. CredTT's contribution is combining them with conditioning as primitive and exploring the consequences for incompleteness and choice. Many results are postulated analogs of proven theorems in simpler settings.

\subsection{What CredTT Rejects}

\begin{enumerate}
    \item \textbf{Implication}: No $A \to B := \negc A \lor B$. Use conditioning instead.
    \item \textbf{Addition}: No primitive $+$. Disjunction derived via De Morgan.
    \item \textbf{Division}: No $P(A|B) = P(A,B)/P(B)$. Use chain rule multiplicatively.
    \item \textbf{MLTT as foundation}: MLTT is the $\{\zero,\one\}$ \emph{limit}, not the starting point.
\end{enumerate}

%=============================================================================
\section{Categorical Semantics}
%=============================================================================

\begin{definition}[Credence-Weighted Category]
A \emph{$\C$-weighted category} has:
\begin{itemize}
    \item Objects (types)
    \item For each $c : \C$, morphisms $\mathrm{Hom}_c(A, B)$
    \item Identity: $\mathrm{id}_A \in \mathrm{Hom}_\one(A, A)$
    \item Composition: $f \in \mathrm{Hom}_c(A, B)$, $g \in \mathrm{Hom}_v(B, C)$ gives $g \circ f \in \mathrm{Hom}_{c \cdot v}(A, C)$
\end{itemize}
\end{definition}

When $\C = {[0,1]}$ with copy $\Delta : A \to A \times A$ and delete $! : A \to \mathbf{1}$, this becomes a \textbf{Markov category} \citep{fritz2020synthetic}.

%=============================================================================
\section{Summary}
%=============================================================================

\textbf{Credence algebra}: $(\zero, \one, \cdot, \negc, \leq)$ --- Product De Morgan

\textbf{Derived}: $c \vee v = \negc(\negc c \cdot \negc v)$ --- De Morgan OR

\textbf{Judgments}: $\Gamma \vdash a : A \at{c}$

\textbf{Key rule}: Credences multiply in elimination

\textbf{Conditioning}: Chain rule $P(A,B) = P(A|B) \cdot P(B)$, no division

\textbf{MLTT}: The $\C = \{\zero, \one\}$ case

\textbf{Probability}: The $\C = {[0,1]}$ case

\textbf{Novelty}: De Morgan algebra + type theory + conditioning

%=============================================================================
\section{Neighbourhood Semantics}
%=============================================================================

We replace truth values with \emph{stability under dynamics}. This is not metaphorical: each proof step $s$ induces a monotone operator $T_s(c) = c \cdot s$, and stability is about \emph{fixed points} of these operators.

\subsection{Proofs as Monotone Dynamical Systems}

The key insight: proofs are not static objects but \emph{dynamical systems} on the credence algebra.

\begin{definition}[Proof Operator]
Each proof step at credence $s$ induces the operator $T_s : \C \to \C$ where $T_s(c) = c \cdot s$.
\end{definition}

Since multiplication is monotone, $T_s$ is a monotone operator. Stability is about fixed points:

\begin{definition}[Post-Fixed Point]
$c$ is a \emph{post-fixed point} under step $s$ (written $c \leq c \cdot s$) iff applying the proof step does not reduce credence.
\end{definition}

\begin{definition}[Invariant]
$c$ is \emph{invariant} under step $s$ (written $c = c \cdot s$) iff it is an exact fixed point.
\end{definition}

\begin{definition}[Degenerating]
$c$ \emph{degenerates} under step $s$ iff $\inf_n (c \cdot s^n) = \zero$.
\end{definition}

\begin{definition}[Idempotent]
$c$ is \emph{idempotent} iff $c \cdot c = c$. Such elements are self-stable and may be \emph{interior} ($\zero < c < \one$).
\end{definition}

\textbf{Critical}: We do NOT assume Archimedeanicity. In non-Archimedean algebras, $s < \one$ does NOT imply $s^n \to \zero$. Interior idempotents can exist, allowing \emph{genuine stability in the interior}.

These are \emph{meta-properties}, like termination, not judgments inside CredTT. NO THRESHOLDS. NO ARBITRARY CONSTANTS. Pure order-theoretic dynamics.

\subsection{Classical Techniques as Stability Theorems}

Every classical proof technique has a CredTT analog expressed in terms of stability. The table below shows all 20 classical techniques with their CredTT translations:

\begin{center}
\small
\begin{tabular}{|c|l|l|l|}
\hline
\textbf{\#} & \textbf{Classical Technique} & \textbf{CredTT Analog} & \textbf{Credence Behavior} \\
\hline
1 & Direct proof & Track credence bounds & $c$ accumulates via $\cdot$ \\
2 & Proof by cases & Context refinement & Each branch stable $\Rightarrow$ stable \\
3 & Contraposition & Stability-contraposition & Not exact; $\negc c_2 \leq \negc c_1$ \\
4 & Reductio ad absurdum & Instability reflection & Unstable$(\negc A) \Rightarrow$ Stable$(A)$ \\
5 & Contrapositive & Robust implication & Stability bounds transform \\
6 & Ex falso & Limit admissibility & As $c \to \zero$, constraints vanish \\
7 & Vacuous truth & Degenerate conditioning & $c \approx \zero \Rightarrow$ unconstrained \\
8 & Deduction theorem & Function abstraction & $\Gamma, x:A \vdash b:B \at{c}$ preserves $c$ \\
9 & Modus ponens & Application rule & $c_1 \cdot c_2$ (multiply!) \\
10 & Syllogism & Function composition & $c_2 \cdot c_1$ (accumulates) \\
11 & Universal generalization & $\Pi$-introduction & $\lambda x.b : \Pi x.B \at{\inf_x c(x)}$ \\
12 & Existential intro/elim & $\Sigma$-types & Credence multiplies in elim \\
13 & Natural induction & Same + stability of step & May degrade unless contractive \\
14 & Strong induction & Bound over premise set & Multiple uses multiply \\
15 & Exhaustion & Finite case + aggregation & Stable if each case stable \\
16 & Construction & Existence + certificate & Credence = witness $\cdot$ property \\
17 & Refutation/resolution & Degeneracy detection & Drive bound to $\zero$ \\
18 & Equational rewriting & J-elimination & Preserves stability if Id stable \\
19 & Proof by analogy & Low-credence morphism & Visible weakness at $c < \one$ \\
20 & Structural rules & Conditional theorems & See \S\ref{sec:structural-credtt} \\
\hline
\end{tabular}
\end{center}

\paragraph{Concrete Example: Syllogism (Technique \#10)}
Consider the classical syllogism ``All men are mortal; Socrates is a man; therefore Socrates is mortal.''
\begin{itemize}
    \item \textbf{Classical}: Major premise + minor premise $\Rightarrow$ definite conclusion.
    \item \textbf{CredTT}: Major @ $c_1$ + Minor @ $c_2$ $\Rightarrow$ Conclusion @ $c_1 \cdot c_2$.
\end{itemize}
If $c_1 = 0.95$ and $c_2 = 0.90$, the conclusion has credence $0.855$---\emph{not} $0.95$ or $0.90$!

\paragraph{Concrete Example: Induction (Technique \#13)}
Induction is valid iff $c$ is a \emph{post-fixed point} of the step operator:
\begin{itemize}
    \item \textbf{Classical}: step at credence $\one$, so $c \leq c \cdot \one = c$ always holds
    \item \textbf{Interior}: if $c$ is idempotent ($c \cdot c = c$), induction at $c$ is valid!
    \item \textbf{Degenerating}: if $c < c \cdot s$ fails, the proof degrades over iterations
\end{itemize}
The formulation $P(n) \at{c_b \cdot c_s^n}$ assumes Archimedean credences. In non-Archimedean algebras with interior idempotents, induction can stabilize at interior credences.

\paragraph{Concrete Example: Ex Falso (Technique \#6)}
The graded spectrum of ex falso:
\begin{itemize}
    \item $c = \one$: Fully constrained (normal logic)
    \item $c = 0.5$: Weakly constrained
    \item $c = 0.01$: Almost unconstrained
    \item $c = \zero$: Completely unconstrained (classical ex falso)
\end{itemize}
The chain rule $P(A,B) = P(A|B) \cdot P(B)$ is trivially satisfied when $P(B) = \zero$, making $P(A|B)$ unconstrained.

\subsection{The $\{\zero,\one\}$ Collapse Theorem}

\begin{theorem}[Collapse]
When credences are restricted to $\{\zero,\one\}$:
\begin{enumerate}
    \item Neighbourhoods become singletons
    \item $\mathrm{Stable}_\one(\one) = \mathrm{true}$, $\mathrm{Unstable}_\zero(\zero) = \mathrm{true}$
    \item CredTT judgments correspond exactly to MLTT judgments
    \item No interior points exist (every credence is extreme)
\end{enumerate}
\end{theorem}

\begin{proof}[Proof sketch]
In Boolean algebra, there are no intermediate points between $\zero$ and $\one$, so $\epsilon$-neighbourhoods collapse to their centers. The stability classification becomes exhaustive and trivial: $\one$ is stable, $\zero$ is unstable.
\end{proof}

This explains why classical logic appears to have exact truth values: it's the degenerate limit where neighbourhoods collapse to points. CredTT reveals the underlying continuous structure that classical logic hides.

%=============================================================================
\section{Proof Principles in CredTT}
%=============================================================================

\subsection{Object-Level vs Meta-Level}

CredTT distinguishes two kinds of reasoning:
\begin{itemize}
    \item \textbf{Object-level}: Produces judgments $\Gamma \vdash a : A \at{c}$
    \item \textbf{Meta-level}: Reasons about stability regions
\end{itemize}

Object-level reasoning uses the type rules. Meta-level reasoning proves things like ``if both premises are stable, the conclusion is stable.''

\subsection{Recovered Classical Patterns}

Classical proof patterns are recovered as:
\begin{itemize}
    \item \textbf{Function transformers}: Deduction, modus ponens, transitivity preserve credence via multiplication
    \item \textbf{Stability reflection}: Reductio works via complement: unstable $\negc A$ implies stable $A$
    \item \textbf{Limit admissibility}: Ex falso is the $c = \zero$ limit where multiplication annihilates
    \item \textbf{Controlled structural rules}: Weakening/contraction are conditional on credence
\end{itemize}

\subsection{New CredTT-Native Techniques}

CredTT enables 8 proof techniques with \emph{no classical analog}. These are the core differentiators:

\begin{center}
\small
\begin{tabular}{|c|l|l|}
\hline
\textbf{\#} & \textbf{Technique} & \textbf{Description} \\
\hline
21 & Stability proofs & Prove ``$A$ is Stable$_\one$'' or ``$B$ is Unstable$_\zero$'' \\
22 & Credence bounds as invariants & Maintain $c(A) \geq l$ through computation \\
23 & Continuity/robustness lemmas & Small input degradation $\Rightarrow$ small output degradation \\
24 & Degeneracy analysis & Quantify how fast $c \to \zero$ (diagnosis) \\
25 & Contractivity arguments & Prove step is non-degrading ($c_{out} \geq c_{in}$) \\
26 & Proof factoring & Isolate stable core from unstable fringe \\
27 & Dual proofs & Squeeze: $l \leq c(A) \leq u$ from both sides \\
28 & Limit theorems & Asymptotic credence behavior as $n \to \infty$ \\
\hline
\end{tabular}
\end{center}

\paragraph{Concrete Example: Stability Proofs (Technique \#21)}
Classical logic cannot express ``$2+2=4$ is \emph{robustly} true.'' CredTT can:
\begin{itemize}
    \item ``$2+2=4$'' has credence $\one$, which is Stable$_\one$.
    \item ``$\sqrt{2}$ is rational'' has credence $\zero$ (after proof), which is Unstable$_\zero$.
    \item The G\"odel sentence has credence $\tfrac{1}{2}$, which is Interior (neither stable nor unstable).
\end{itemize}

\paragraph{Concrete Example: Dual Proofs (Technique \#27)}
The G\"odel sentence $G$ satisfies:
\begin{itemize}
    \item Lower bound: $c(G) \geq \tfrac{1}{2}$ (from self-reference)
    \item Upper bound: $c(G) \leq \tfrac{1}{2}$ (from negation self-reference)
    \item Squeezed: $c(G) = \tfrac{1}{2}$ exactly
\end{itemize}
This is the \emph{Interior region} that classical logic cannot express. The fixpoint $c = \negc c$ has solution $c = \tfrac{1}{2}$.

\paragraph{Concrete Example: Degeneracy Analysis (Technique \#24)}
In a complex proof with 100 steps:
\begin{itemize}
    \item Steps 1--50: $c$ remains at 0.95
    \item Step 51: $c$ drops to 0.5 (\textbf{problematic step!})
    \item Steps 52--100: $c$ gradually reaches 0.1
\end{itemize}
Degeneracy analysis pinpoints Step 51 as responsible for 80\% of credence loss---like profiling for proofs.

\subsection{Structural Rules}
\label{sec:structural-credtt}

Structural rules have subtle behavior in CredTT---this is Technique \#20 from the classical list:

\begin{itemize}
    \item \textbf{Exchange}: Holds unconditionally. Reordering $\Gamma, A, B, \Delta$ to $\Gamma, B, A, \Delta$ preserves credence. Order doesn't matter.

    \item \textbf{Weakening}: \emph{Conditional} on added assumptions having $c = \one$. If $\Gamma \vdash t : A \at{c}$, then $\Gamma, x:B \at{\one} \vdash t : A \at{c}$. But adding $x:B \at{d}$ with $d < \one$ may degrade conclusions.

    \item \textbf{Contraction}: Requires explicit duplication permission. Using assumption $A \at{c}$ twice might cost $c^2$, not $c$. In resource-sensitive settings, this tracks consumption.
\end{itemize}

Classical logic has free structural rules---this \emph{hides} important credence costs. CredTT makes these costs \emph{explicit}.

%=============================================================================
\section{Related Work: Neighbourhood and Continuous Semantics}
%=============================================================================

\subsection{Primitive Conditional Probability}

The tradition of taking conditional probability as primitive (rather than derived via the ratio formula) traces to \textbf{Popper functions} and \textbf{R\'enyi conditional probability}. Stalnaker \citep{stalnaker1970probability} and Pruss \citep{pruss2013popper} develop this approach philosophically. CredTT adopts this stance but internalizes it into dependent type theory.

\subsection{Continuous Logic and Approximate Satisfaction}

\textbf{Continuous logic} for metric structures uses truth values in $[0,1]$ with robustness/continuity as key notions \citep{benyaacov2008continuous}. This resonates with our ``truth as stability'' view, but continuous logic is model-theoretic rather than proof-theoretic.

\subsection{Neighbourhood Semantics}

Many-valued generalizations of neighbourhood semantics exist \citep{hansoul2020neighbourhood}, compatible with our stability regions, but not in a dependent type setting.

\subsection{What CredTT Unifies}

CredTT sits at the intersection of:
\begin{itemize}
    \item Popper-style primitive conditioning
    \item Continuous/approximate logic robustness
    \item Graded type theory judgment structure
\end{itemize}

But differs by: (i) conditioning-first semantics, (ii) $\bot$/$\top$ as neighbourhood phenomena, (iii) MLTT as the $\{\zero,\one\}$ collapse.

%=============================================================================
\section{Graded Incompleteness}
%=============================================================================

G\"odel's incompleteness theorems have a natural formulation in CredTT. Rather than forcing undecidability, self-referential sentences attain determinate intermediate credences.

\subsection{Graded Provability Predicate}

\begin{definition}[Graded Provability]
For a theory $T$ and proposition $\varphi$, define:
\[
\mathsf{Prov}_c(\ulcorner \varphi \urcorner) :\equiv \text{there exists a derivation of } \varphi \text{ at credence } c
\]
The \emph{provability degree} is:
\[
|\varphi|_T := \sup\{c \mid T \vdash \varphi \at{c}\}
\]
\end{definition}

This follows Pavelka \citep{pavelka1979fuzzy} and Hajek \citep{hajek1998metamathematics}, who developed graded provability for fuzzy logics.

\subsection{The Diagonal Lemma}

The diagonal lemma is purely syntactic and holds for any sufficiently expressive language:

\begin{lemma}[Diagonal Lemma]
For any predicate $P(x)$, there exists a sentence $G$ such that:
\[
T \vdash G \leftrightarrow P(\ulcorner G \urcorner)
\]
\end{lemma}

In CredTT, this becomes: for any credence-valued predicate $P$, there exists $G$ such that the credence of $G$ equals $P(\ulcorner G \urcorner)$.

\subsection{The G\"odel Sentence in CredTT}

\begin{definition}[G\"odel Sentence]
Let $G$ be the sentence asserting ``$G$ is not provable at credence $\one$'':
\[
G \leftrightarrow \neg \mathsf{Prov}_\one(\ulcorner G \urcorner)
\]
\end{definition}

\begin{theorem}[Graded Incompleteness]
If $G$ has credence $c$ in CredTT, then $c$ satisfies $c = \negc c$, hence $c = \frac{1}{2}$.
\end{theorem}

\begin{proof}
By the definition of $G$: if $G$ has credence $c$, then $G$'s content (that $G$ is not provable at credence $\one$) should also have credence $c$. But $G$'s claim is essentially the negation of provability. In the fixed-point analysis:
\begin{itemize}
    \item If $c = \one$, then $G$ asserts $\neg\mathsf{Prov}_\one(G)$, but $G$ itself is provable at $\one$---contradiction.
    \item If $c = \zero$, then $\neg G$ has credence $\one$, meaning $G$ is provable at $\one$---but we assumed $c = \zero$.
\end{itemize}
The consistent resolution: $c = \negc c$. With $\negc c = 1 - c$:
\[
c = 1 - c \quad \Rightarrow \quad c = \frac{1}{2}
\]
\end{proof}

\begin{remark}[Comparison with Classical Incompleteness]
\begin{center}
\begin{tabular}{l|l|l}
& \textbf{Classical} & \textbf{CredTT} \\
\hline
$G$'s status & Undecidable & Credence $\frac{1}{2}$ \\
$c = \negc c$ in $\{\zero, \one\}$ & No solution & N/A \\
$c = \negc c$ in $[0,1]$ & N/A & $c = \frac{1}{2}$ \\
Information & Binary gap & Quantified uncertainty \\
\end{tabular}
\end{center}
\end{remark}

The key insight: classical logic's restriction to $\{\zero, \one\}$ forces $c = \negc c$ to have no solution, manifesting as undecidability. CredTT's continuous credences allow the fixed point to exist at $\frac{1}{2}$.

\subsection{Literature Context}

The fixed point $c = \negc c \Rightarrow c = \frac{1}{2}$ is well-studied in fuzzy logic:

\begin{itemize}
    \item Hajek, Paris, Shepherdson, ``The Liar Paradox and Fuzzy Logic'' (JSL 2000) \citep{hajek2000liar}
    \item Restall on Lukasiewicz logic and self-reference
\end{itemize}

CredTT's contribution is embedding this in dependent type theory with graded judgments.

%=============================================================================
\section{Graded Self-Consistency}
%=============================================================================

G\"odel's Second Incompleteness Theorem states that a consistent theory extending PA cannot prove its own consistency. CredTT offers a graded alternative.

\subsection{The Classical Obstacle}

\begin{theorem}[G\"odel's Second, Classical]
If $T \supseteq PA$ and $T$ is consistent, then $T \not\vdash \mathsf{Con}(T)$.
\end{theorem}

Here $\mathsf{Con}(T) := \neg\mathsf{Prov}_T(\ulcorner \bot \urcorner)$.

\subsection{Graded Consistency}

\begin{definition}[Graded Consistency]
\[
\mathsf{Con}_c(T) :\equiv \text{``$\bot$ is not provable at credence $\one$''} \at{c}
\]
\end{definition}

\begin{theorem}[Graded Second Incompleteness]
CredTT can prove $\mathsf{Con}_c(\text{CredTT})$ for some $c < \one$, but not for $c = \one$.
\end{theorem}

\begin{proof}[Proof sketch]
The G\"odel-2 argument requires:
\begin{enumerate}
    \item If $T \vdash \mathsf{Con}(T) \at{\one}$, then by formalized soundness, $T$ proves ``if $T$ proves $\bot$, then $\bot$''
    \item Combined with the G\"odel sentence, this yields $T \vdash G \at{\one}$
    \item But $G$ asserts its own unprovability at $\one$---contradiction
\end{enumerate}

At credence $c < \one$, step 1 fails: we only have $\mathsf{Con}(T) \at{c}$, which doesn't yield the sharp consequence needed. The self-referential loop stabilizes at a credence below $\one$.
\end{proof}

\subsection{Philosophical Interpretation}

\begin{center}
\begin{tabular}{l|l}
\textbf{Classical} & \textbf{CredTT} \\
\hline
Cannot prove own consistency & Can prove at $c < \one$ \\
``Unknowable'' & ``Confident but not certain'' \\
Binary epistemic gap & Quantified confidence level \\
\end{tabular}
\end{center}

This is philosophically honest: we \emph{believe} our system is consistent, but absolute certainty is impossible (G\"odel). CredTT quantifies this epistemic state rather than leaving it as a binary gap.

\subsection{Connection to LFIs}

Logics of Formal Inconsistency \citep{carnielli2016paraconsistent} use an object-level consistency operator $\circ A$ (``$A$ is consistent''). The ``gentle explosion'' principle:
\[
A, \neg A, \circ A \vdash B \quad \text{but} \quad A, \neg A \not\vdash B
\]

CredTT achieves similar effects through credence: contradiction at low credence doesn't explode to arbitrary conclusions because credence multiplication preserves the low credence.

\subsection{Self-Hosting via Credence-Valued Provability}

CredTT can describe itself internally. The key design choice:

\begin{definition}[Internal Provability]
$\mathsf{Prov}(\varphi) : \C$ maps formulas (via G\"odel coding) to credences, representing the ``degree of provability.''
\end{definition}

This is \emph{credence-valued}, not threshold-based. No arbitrary constants.

\textbf{What this enables:}
\begin{itemize}
    \item Syntax encoding: CredTT represents its own terms, types, derivations
    \item Meta-theorems: Internal reasoning about credence propagation
    \item Self-reference: The G\"odel sentence $G \simeq \negc\mathsf{Prov}(G)$ has solution $\mathsf{Prov}(G) = \frac{1}{2}$
\end{itemize}

\textbf{Reflection principle (non-expansive):}
\[
\mathsf{Prov}(\varphi) \leq \mathsf{Sem}(\varphi)
\]
What we prove internally never exceeds semantic truth. For meta-level reasoning, a discount factor prevents L\"ob-style self-certification loops.

\textbf{Graded inconsistency:} If $P \at{c}$ and $\negc P \at{d}$ with $c,d > \zero$, the inconsistency credence is $c \cdot d$. This propagates but doesn't explode. Ex falso becomes limit admissibility: as $c \to \zero$, constraints vanish.

%=============================================================================
\section{Graded Choice}
%=============================================================================

The Axiom of Choice has a natural graded formulation in CredTT.

\subsection{Classical AC and Its Problems}

\begin{definition}[Axiom of Choice]
For any family of non-empty sets $\{A_i\}_{i \in I}$, there exists a choice function $f : I \to \bigcup_i A_i$ with $f(i) \in A_i$ for all $i$.
\end{definition}

AC leads to non-constructive results like the Banach-Tarski paradox: a ball can be decomposed and reassembled into two balls of the same size.

\subsection{Graded Choice}

\begin{definition}[Graded AC]
\[
\mathsf{AC}_c : \left(\prod_{i:I} \|A_i\|_{c_i}\right) \to \left\|\prod_{i:I} A_i\right\|_c
\]
where $c$ depends on $I$ and the $c_i$.
\end{definition}

\begin{theorem}[Choice Credence Scaling]
\begin{enumerate}
    \item \textbf{Finite choice}: $c = \one$ (fully constructible)
    \item \textbf{Countable choice}: $c = \prod_{n=1}^\infty c_n$ (may be $< \one$)
    \item \textbf{Uncountable choice}: $c \ll \one$
\end{enumerate}
\end{theorem}

\begin{proof}[Proof sketch]
For finite $I$, we can construct the choice function explicitly---each step preserves credence $\one$, and finitely many steps preserve the product.

For countable $I$, the infinite product $\prod_{n=1}^\infty c_n$ converges but may be strictly less than $\one$ (e.g., if $c_n = 1 - \epsilon/2^n$, then $c < 1$).

For uncountable $I$, the product over uncountably many factors $< \one$ collapses toward $\zero$.
\end{proof}

\subsection{Banach-Tarski at Credence $< \one$}

\begin{theorem}[Graded Banach-Tarski]
The Banach-Tarski decomposition exists at some credence $c_{BT} < \one$.
\end{theorem}

\begin{remark}[Physical Interpretation]
Credence $< \one$ means ``exists mathematically but not physically realizable.'' Banach-Tarski requires infinitely many non-constructive choices; each contributes a factor $< \one$; the product is strictly $< \one$. This correctly predicts that you cannot actually perform the decomposition.
\end{remark}

\subsection{Comparison}

\begin{center}
\begin{tabular}{l|l|l|l}
& \textbf{Constructive} & \textbf{Classical} & \textbf{CredTT} \\
\hline
AC status & Fails & Holds (credence $\one$) & Holds at credence $c$ \\
BT paradox & Doesn't exist & Exists & Exists at $c < \one$ \\
Information & Binary rejection & Binary acceptance & Graded \\
\end{tabular}
\end{center}

CredTT provides more information than both: constructive math says ``no choice function''; classical math says ``choice function exists''; CredTT says ``choice function exists at credence $c$, quantifying constructibility.''

%=============================================================================
\section{Implementation Status}
%=============================================================================

\subsection{What Is Implemented}

\begin{enumerate}
    \item \textbf{Credence algebra} (\texttt{credtt-impl/src/credence.ml}): Product De Morgan structure with:
    \begin{itemize}
        \item Simplification of credence expressions
        \item Rational arithmetic for exact computation
        \item Fixpoint solver for $c = \negc c \Rightarrow c = \frac{1}{2}$
        \item Unification-based constraint accumulation
    \end{itemize}
    \item \textbf{Bidirectional type checker} (\texttt{credtt-impl/src/check.ml}): Prototype implementing:
    \begin{itemize}
        \item Credence multiplication in elimination rules
        \item Type inference with credence tracking
        \item Credence comparison via $\leq$
    \end{itemize}
    \item \textbf{Proof checker} (\texttt{credtt-impl/src/proof.ml}): Demonstration system for:
    \begin{itemize}
        \item Postulates, derivations, contradictions
        \item Credence forcing from contradictions
        \item Negation fixpoint solving
    \end{itemize}
    \item \textbf{Agda formalization} (\texttt{agda/CredTT/}): Syntax, judgments, MLTT embedding theorem
\end{enumerate}

\subsection{What Is NOT Implemented}

\begin{enumerate}
    \item \textbf{Credence inference}: The checker does not infer credences from term structure. Credences must be explicitly annotated or default to $\one$. The \texttt{infer\_ctx} in \texttt{credence.ml} accumulates constraints but \texttt{solve\_inference} handles only trivial cases.
    \item \textbf{Constraint propagation}: The Pi-Elim rule should generate $c_{result} = c_f \cdot c_a$, but constraints are not propagated bidirectionally.
    \item \textbf{Dependent credence instantiation}: Credences depending on bound variables (e.g., $c(x)$ in $(x : A) \to B \at{c(x)}$) are parsed but not fully instantiated during application.
    \item \textbf{General proof checking}: The proof checker (\texttt{proof.ml}) demonstrates fixed examples (e.g., $\sqrt{2}$ irrationality). It does not parse arbitrary proof files.
\end{enumerate}

\subsection{Prototype vs Production}

The implementation is a \textbf{prototype} demonstrating the core ideas:
\begin{itemize}
    \item Variables are always assigned credence $\one$ (line 24 of \texttt{check.ml})
    \item Non-trivial fixpoints like $c = c \cdot c$ choose $1$ arbitrarily
    \item The $\sqrt{2}$ proof runs a hardcoded list of declarations, not a parser
\end{itemize}

A production implementation would require bidirectional credence inference, general fixpoint solving, and a proper parser for proof files.

%=============================================================================
\section{Open Problems}
%=============================================================================

\begin{enumerate}
    \item \textbf{Credence inference}: Can credences be inferred from term structure?
    \item \textbf{Decidability}: Type checking decidable if $\C$ decidable---proof needed
    \item \textbf{Normalization}: Should follow from MLTT---formal proof needed
    \item \textbf{Dependent credences}: $c(x)$ varying over $x$ in $(x : A) \to B \at{c(x)}$
    \item \textbf{Continuous types}: Extension to measurable spaces
    \item \textbf{Full Agda proofs}: Some metatheorems have holes (\texttt{--allow-unsolved-metas})
    \item \textbf{Completeness}: Pavelka-style completeness for CredTT?
\end{enumerate}

\bibliographystyle{plainnat}
\bibliography{credtt}

\end{document}
